% ============================================================
%  MOSHO — Ablation Study  (Reorganised)
%  Metrics : IGD+, nHV, Spacing, Exec-Time
%  Units   : U01–U07 (base) · U08, U10, U11 (advanced)
% ============================================================
\documentclass[11pt,a4paper]{article}

% ---------- packages ----------
\usepackage[T1]{fontenc}
\usepackage[utf8]{inputenc}
\usepackage[margin=2.5cm]{geometry}
\usepackage{booktabs}
\usepackage{multirow}
\usepackage{array}
\usepackage{xcolor}
\usepackage{colortbl}
\usepackage{amsmath,amssymb}
\usepackage{graphicx}
\usepackage{caption}
\usepackage{subcaption}
\usepackage{hyperref}
\usepackage{enumitem}
\usepackage{tabularx}
\usepackage{makecell}
\usepackage{rotating}
\usepackage{pdflscape}
\usepackage{float}
\usepackage{tikz}
\usepackage{pgfplots}
\pgfplotsset{compat=1.18}
\usepackage{pgfplotstable}
\usepackage{siunitx}
\usepackage{boldline}
\usepackage{longtable}
\usepackage{pifont}

% ---------- colours ----------
\definecolor{critred}{HTML}{C0392B}
\definecolor{imptorange}{HTML}{E67E22}
\definecolor{modblue}{HTML}{2980B9}
\definecolor{marggreen}{HTML}{27AE60}
\definecolor{rowgray}{HTML}{F2F2F2}
\definecolor{headerblue}{HTML}{1A3A5C}
\definecolor{bestgreen}{HTML}{D5F5E3}
\definecolor{worstred}{HTML}{FADBD8}
\definecolor{pipelinered}{HTML}{E74C3C}
\definecolor{pipelineorange}{HTML}{E67E22}
\definecolor{pipelineyellow}{HTML}{F39C12}
\definecolor{pipelinegreen}{HTML}{27AE60}

% ---------- helpers ----------
\newcommand{\up}{\textcolor{critred}{$\uparrow$}}
\newcommand{\dn}{\textcolor{marggreen}{$\downarrow$}}
\newcommand{\best}[1]{\cellcolor{bestgreen}\textbf{#1}}
\newcommand{\worst}[1]{\cellcolor{worstred}{#1}}
\newcommand{\mosho}{\textsc{Mosho}}
\newcommand{\mowso}{\textsc{Mowso}}
\newcommand{\ablof}[1]{\texttt{ABL-#1}}

% ---------- contribution-tier labels ----------
\newcommand{\critical}[1]{\textbf{\textcolor{critred}{#1}}}
\newcommand{\important}[1]{\textbf{\textcolor{imptorange}{#1}}}
\newcommand{\secondary}[1]{\textbf{\textcolor{modblue}{#1}}}
\newcommand{\neutral}[1]{\textcolor{gray}{#1}}

% ---------- pipeline step labels ----------
\newcommand{\stepmand}[2]{%
  \noindent\textcolor{#1}{\rule{3pt}{10pt}}\;%
  \textbf{#2}\quad\textit{(Mandatory)}}
\newcommand{\stepopt}[1]{%
  \noindent\textcolor{pipelinegreen}{\rule{3pt}{10pt}}\;%
  \textbf{#1}\quad\textit{(Optional --- recommended)}}

\hypersetup{
  colorlinks=true,
  linkcolor=headerblue,
  citecolor=headerblue,
  urlcolor=headerblue,
  pdftitle={MOSHO Ablation Study},
  pdfauthor={},
}

% ============================================================
\begin{document}

% ------------------------------------------------------------ title page
\begin{titlepage}
  \centering
  \vspace*{2cm}
  {\color{headerblue}\rule{\linewidth}{1.5pt}}\\[6pt]
  {\LARGE\bfseries\color{headerblue}
    Ablation Study\\[4pt]
    \mosho{} --- Multi-Objective Shark Hunting Optimization\\[4pt]
    for Neural Architecture Search}\\[6pt]
  {\color{headerblue}\rule{\linewidth}{1.5pt}}\\[1.5cm]
  {\large\itshape Rigorous evaluation of the individual and combined
   contribution\\of the ten algorithmic units}\\[2cm]
 
  {\small\color{gray} Draft --- filled from \texttt{experiments/results\_ablation\_study}}
\end{titlepage}

\tableofcontents
\clearpage

% ============================================================
\section{Introduction}
\label{sec:intro}
% ============================================================

\subsection{Overview and Objectives}
\label{sec:intro_overview}

This document describes the ablation study protocol for \mosho{}, a
multi-objective meta-heuristic designed for Neural Architecture Search
(NAS). \mosho{} incorporates \textbf{ten algorithmic units} beyond its
predecessor \mowso{}: seven constituting the \emph{base layer} (U01--U07)
and three constituting the \emph{advanced layer} (U08, U10, U11).

The study isolates and quantifies the individual contribution of each unit
through controlled one-at-a-time (OAT) ablations, group ablations, and
rigorous statistical validation.  Four complementary metrics are used
throughout: IGD$^+$, normalised hypervolume (nHV), Spacing, and execution
time.  The central questions addressed are:

\begin{enumerate}[leftmargin=2em, itemsep=3pt]
  \item Does removing unit $U_i$ significantly degrade performance?
  \item What is the \emph{magnitude} of each unit's individual contribution?
  \item Do any pairs or groups of units interact super-additively?
 
\end{enumerate}

\subsection{Description of the Algorithmic Units}
\label{sec:intro_units}


\begin{longtable}{>{\bfseries\ttfamily}p{1.0cm} p{2.8cm} p{5.4cm} p{4.4cm}}
\toprule
\textbf{Code} & \textbf{Name} &
\textbf{Description in \mosho{}} & \textbf{Ablation replacement} \\
\midrule
\endfirsthead
\multicolumn{4}{l}{\small\itshape (continued)}\\
\toprule
\textbf{Code} & \textbf{Name} &
\textbf{Description in \mosho{}} & \textbf{Ablation replacement} \\
\midrule
\endhead
\bottomrule
\endlastfoot

U01 & Discrete kernels &
  Six integer operators (Patrol, Scent, Circle, Burst, Crossover, Scout)
  producing valid architecture vectors directly, without rounding. &
  Continuous update $\mathbf{x}(t{+}1) = \mathbf{x}(t)+\mathbf{v}(t{+}1)$
  followed by quantisation $Q(\cdot)$ (\mowso{} behaviour). \\[4pt]

U02 & Uniform crossover &
  Gene-by-gene mixing with a partner ($p{=}0.5$ per gene) followed by
  single-gene mutation ($p{=}0.3$). &
  No crossover --- each agent moves independently (\mowso{} behaviour). \\[4pt]

U03 & Deferred truncation &
  Dominance check $O(A)$ inside the loop; one crowding-distance pass
  $O(A\log A)$ at end of iteration. &
  True-distance recomputed after every insertion; complexity
  $O(N{\cdot}A^3)$ per iteration (\mowso{} behaviour). \\[4pt]

U04 & Explicit acceptance &
  Three-case rule: (1)~dominance $\Rightarrow$ always accepted;
  (2)~archive improvement $\Rightarrow$ accepted;
  (3)~mutual non-dominance $\Rightarrow$ $p{=}0.4{\times}(1{-}0.7p_t)$. &
  Unconditional acceptance of any candidate (\mowso{} behaviour). \\[4pt]

U05 & Energy + restart &
  Energy $e_i\!\in\![0,e_{\max}]$ updated each step
  ($\beta_1{=}0.20$ if archive improved, $\beta_2{=}0.05$ otherwise,
  $\delta{=}0.0968$ if rejected); restart when $e_i<e_{\min}{=}0.508$. &
  No restart --- stagnant agents remain blocked indefinitely
  (\mowso{} behaviour). \\[4pt]

U06 & Dual adaptive signal &
  $\boldsymbol{\pi}(t)=\mathrm{Normalise}(\boldsymbol{\pi}_0\odot
  g(r_t,p_t))$ adapted according to improvement rate $r_t$
  (window $H{=}20$) and progress $p_t{=}\text{evals}/B$. &
  Fixed exponential decays for $p_1,p_2$, independent of runtime state
  (\mowso{} behaviour). \\[4pt]

U07 & Active archive &
  Archive drives: (i)~operator probabilities via $r_t$,
  (ii)~restarts by recombination, (iii)~acceptance decisions (Case~2). &
  Passive archive: stores non-dominated solutions and provides
  a single leader (max true distance) (\mowso{} behaviour). \\[4pt]


U08 & K-means archive &
  K-means on normalised objective vectors
  ($K=\max(2,\mathrm{round}(\sqrt{|A|}))$ when $|A|\ge 20$).
  Cluster $C_k$ selected with $P(C_k)\propto 1/|C_k|$,
  then uniform draw within $C_k$. &
  Uniform sampling from archive (each solution equally probable). \\[4pt]

U10 & Extreme solutions &
  Registry $\mathbf{x}_{\mathrm{ext}}=\arg\min/\max f_i$ per objective.
  Excluded from truncation; local refinement budget allocated. &
  No protection --- boundary solutions may be removed by standard
  crowding-distance truncation. \\[4pt]

U11 & Mono-objective fallback &
  Latency degeneracy detection: if the archive latency range is below
  $\varepsilon{=}10^{-3}$ for three checks, switch acceptance toward the
  accuracy objective. &
  No detection --- algorithm continues in multi-objective mode even when
  one objective collapses. \\
\end{longtable}




% ============================================================
\section{Study Design}
\label{sec:design}
% ============================================================

\subsection{Ablation Variants}
\label{sec:design_variants}

The configured ablation variants are summarised in Table~\ref{tab:variants}.
A~\checkmark\ indicates the unit is \emph{active}; a cross~$\times$ means
it is replaced by its \mowso{}-equivalent reference.  The matrix is
organised in three blocks: (i)~reference configurations, (ii)~OAT
ablations (one unit disabled at a time), and (iii)~group ablations
(multiple units disabled simultaneously).  The current execution includes
full \mosho{}, MOSHO-Base, all OAT ablations, and all group ablations;
\mowso{} is shown as a reference design but was not run in
\texttt{results\_ablation\_study}.

\begin{table}[H]
\centering
\caption{Complete binary matrix of ablation variants
  (\checkmark~= active, $\times$~= disabled).}
\label{tab:variants}
\setlength{\tabcolsep}{4pt}
\renewcommand{\arraystretch}{1.15}
\small
\begin{tabular}{l *{10}{c}}
\toprule
\textbf{Variant} &
  \textbf{U01} & \textbf{U02} & \textbf{U03} & \textbf{U04} &
  \textbf{U05} & \textbf{U06} & \textbf{U07} &
  \textbf{U08} & \textbf{U10} & \textbf{U11} \\
\midrule
\multicolumn{11}{l}{\textit{Reference configurations}}\\[2pt]
\rowcolor{rowgray}
\mosho{} (full) &\checkmark&\checkmark&\checkmark&\checkmark&\checkmark
                &\checkmark&\checkmark&\checkmark&\checkmark&\checkmark\\
MOSHO-Base      &\checkmark&\checkmark&\checkmark&\checkmark&\checkmark
                &\checkmark&\checkmark&$\times$  &$\times$  &$\times$  \\
\mowso{}        &$\times$  &$\times$  &$\times$  &$\times$  &$\times$
                &$\times$  &$\times$  &$\times$  &$\times$  &$\times$  \\
\midrule
\multicolumn{11}{l}{\textit{OAT ablations}}\\[2pt]
\ablof{U01} &$\times$  &\checkmark&\checkmark&\checkmark&\checkmark&\checkmark&\checkmark&\checkmark&\checkmark&\checkmark\\
\ablof{U02} &\checkmark&$\times$  &\checkmark&\checkmark&\checkmark&\checkmark&\checkmark&\checkmark&\checkmark&\checkmark\\
\ablof{U03} &\checkmark&\checkmark&$\times$  &\checkmark&\checkmark&\checkmark&\checkmark&\checkmark&\checkmark&\checkmark\\
\ablof{U04} &\checkmark&\checkmark&\checkmark&$\times$  &\checkmark&\checkmark&\checkmark&\checkmark&\checkmark&\checkmark\\
\ablof{U05} &\checkmark&\checkmark&\checkmark&\checkmark&$\times$  &\checkmark&\checkmark&\checkmark&\checkmark&\checkmark\\
\ablof{U06} &\checkmark&\checkmark&\checkmark&\checkmark&\checkmark&$\times$  &\checkmark&\checkmark&\checkmark&\checkmark\\
\ablof{U07} &\checkmark&\checkmark&\checkmark&\checkmark&\checkmark&\checkmark&$\times$  &\checkmark&\checkmark&\checkmark\\
\ablof{U08} &\checkmark&\checkmark&\checkmark&\checkmark&\checkmark&\checkmark&\checkmark&$\times$  &\checkmark&\checkmark\\
\ablof{U10} &\checkmark&\checkmark&\checkmark&\checkmark&\checkmark&\checkmark&\checkmark&\checkmark&$\times$  &\checkmark\\
\ablof{U11} &\checkmark&\checkmark&\checkmark&\checkmark&\checkmark&\checkmark&\checkmark&\checkmark&\checkmark&$\times$  \\
\midrule
\multicolumn{11}{l}{\textit{Group ablations}}\\[2pt]
G-SEARCH  &$\times$&$\times$&\checkmark&\checkmark&\checkmark&\checkmark&\checkmark&\checkmark&\checkmark&\checkmark\\
G-ADAPT   &\checkmark&\checkmark&\checkmark&\checkmark&$\times$&$\times$&$\times$&\checkmark&\checkmark&\checkmark\\
G-ARCHIVE &\checkmark&\checkmark&\checkmark&\checkmark&\checkmark&\checkmark&\checkmark&$\times$&$\times$&\checkmark\\
G-NOADV   &\checkmark&\checkmark&\checkmark&\checkmark&\checkmark&\checkmark&\checkmark&$\times$&$\times$&$\times$\\
G-NOBASE  &$\times$&$\times$&$\times$&$\times$&$\times$&$\times$&$\times$&\checkmark&\checkmark&\checkmark\\
G-CORE    &\checkmark&$\times$&$\times$&\checkmark&\checkmark&\checkmark&\checkmark&$\times$&$\times$&$\times$\\
\bottomrule
\end{tabular}
\end{table}

\subsection{Experimental Setup}
\label{sec:design_setup}

\subsubsection{Hyperparameter Settings}
\label{sec:design_params}

All variants share the hyperparameter configuration in
Table~\ref{tab:params}.  Seeds are fixed identically across variants to
enable paired statistical comparisons.

\begin{table}[H]
\centering
\caption{Hyperparameter settings shared across all variants.}
\label{tab:params}
\renewcommand{\arraystretch}{1.2}
\begin{tabularx}{\linewidth}{l l X}
\toprule
\textbf{Parameter} & \textbf{Value} & \textbf{Justification} \\
\midrule
Population size $N$           & 50
  & Standard for multi-objective swarms \\
Archive size $M$              & 50
  & Equal to $N$ for fair comparison \\
Iterations                    & 200
  & Computed as $B/N = 10\,000/50$ \\
Evaluation budget $B$         & 10\,000
  & Precision/cost trade-off \\
Repetitions $R$               & 30
  & Minimum for non-parametric tests ($p < 0.05$) \\
Random seeds & $42 + r \times 7,\; r=0,\dots,29$
  & Full reproducibility, identical across variants \\
Datasets & CIFAR-10, CIFAR-100, ImageNet16-120
  & Standard NAS-Bench-201 dataset contexts \\
Devices & edgegpu, edgetpu, eyeriss, fpga, pixel3, raspi4
  & Six hardware-latency contexts per dataset \\
Sliding window $H$            & 20
  & Improvement rate $r_t$ computation (U06) \\
Initial energy $e_0$          & 0.818
  & Initial shark energy in the implementation \\
$\eta$                       & 15.0
  & MOSHO constructor parameter retained for reproducibility \\
$\beta_1$ (energy gain, archive improved) & 0.20
  & Fast gain when archive improves (U05) \\
$\beta_2$ (energy gain, accepted)         & 0.05
  & Slow gain otherwise (U05) \\
$\delta$ (energy decay)       & 0.0968
  & Slow decay to avoid premature restarts \\
$e_{\min}$ (restart threshold)& 0.508
  & Restart trigger (U05) \\
Maximum energy $e_{\max}$     & 1.699
  & Upper energy cap \\
Cluster activation threshold (U08) & $5^6=15\,625 < 100\,000$
  & Cluster sampling is disabled for this NAS-Bench-201 search space \\
K-means minimum archive size (U08) & 20
  & Minimum archive size before cluster sampling can activate \\
K-means refresh / iterations (U08) & 4 / 12
  & Cache refresh period and maximum Lloyd iterations \\
Collapse threshold $\varepsilon$ (U11) & $10^{-3}$
  & Latency degeneracy detection threshold \\
Mono-objective patience / refresh (U11) & 3 / 12
  & Number of checks before activation and refresh interval \\
Post-crossover mutation prob.\ (U02) & 0.30
  & Single-gene mutation after crossover \\
Base operator probabilities & 0.20 / 0.25 / 0.15 / 0.10 / 0.20 / 0.10
  & Patrol / Scent / Circle / Burst / Crossover / Scout \\
Mutual non-dominance acceptance & $0.4(1-0.7p_t)$
  & U04 acceptance probability schedule \\
\bottomrule
\end{tabularx}
\end{table}

\subsubsection{Performance Metrics}
\label{sec:design_metrics}

Four metrics capture complementary facets of Pareto front quality
(Table~\ref{tab:metrics}).
\begin{table}[H]
\centering
\caption{Evaluation metrics and optimisation directions}
\label{tab:metrics}
\renewcommand{\arraystretch}{1.5} % Augmenté un peu pour l'aération

\begin{tabularx}{\linewidth}{@{} l >{\raggedright\arraybackslash}p{3.5cm} >{\raggedright\arraybackslash}X l c @{}}
\toprule
\textbf{Symbol} & \textbf{Full name} & \textbf{Definition / Interpretation} & \textbf{Facet} & \textbf{Dir.} \\
\midrule

$\mathrm{IGD}^+$ & 
Inverted Generational Distance Plus & 
Modified IGD measuring the Euclidean distance from each reference point to the nearest solution in the approximation set, clipped at zero. It penalises poor coverage while remaining Pareto-compliant. & 
Conv. + Cov. & $\downarrow$ \\ \addlinespace

$\mathrm{nHV}$ & 
Normalised Hypervolume & 
Hypervolume dominated by the approximation front, normalised by a reference hypervolume, yielding a score in $[0,1]$. & 
Conv. + Div. & $\uparrow$ \\ \addlinespace

$\mathrm{SP}$ & 
Spacing & 
Standard deviation of distances between consecutive solutions along the approximation front, reflecting distribution uniformity. & 
Uniformity & $\downarrow$ \\ \addlinespace

$T$ & 
Execution Time & 
Wall-clock time (in seconds) per run, used to assess computational overhead (notably for U03 and U08). & 
Efficiency & $\downarrow$ \\

\bottomrule
\end{tabularx}
\end{table}
% ============================================================
\section{Experimental Results}
\label{sec:results}
% ============================================================

All tables report \textbf{median values over 30 paired runs}, after
averaging the 18 context-level measurements within each run.  The values
in this section are computed directly from
\texttt{experiments/results\_ablation\_study}.  The execution contains
full \mosho{}, MOSHO-Base, ten OAT ablations, and six group ablations.
The \mowso{} baseline was not included in this execution root.

\subsection{Median Performance --- OAT Ablations and Baselines}
\label{sec:results_oat}

\begin{table}[H]
\centering
\caption{Median performance --- OAT ablations and reference
  configurations.
  $\downarrow$~lower is better; $\uparrow$~higher is better.}
\label{tab:res_oat}
\renewcommand{\arraystretch}{1.25}
\setlength{\tabcolsep}{6pt}
\small
\begin{tabular}{l
    S[table-format=1.4]
    S[table-format=1.4]
    S[table-format=1.4]
    S[table-format=2.2]}
\toprule
\textbf{Variant} &
  {\textbf{IGD}$^+$ $\downarrow$} &
  {\textbf{nHV} $\uparrow$} &
  {\textbf{SP} $\downarrow$} &
  {\textbf{Time (s)} $\downarrow$} \\
\midrule
\mosho{} (full) & 0.0202 & 0.9984 & 5.3494 & 0.65 \\
MOSHO-Base      & 0.0249 & 0.9983 & 5.4065 & 0.62 \\
\mowso{}        & \multicolumn{4}{c}{not run in \texttt{results\_ablation\_study}} \\
\midrule
\ablof{U01}     & 0.0206 & 0.9981 & 5.3636 & 0.61 \\
\ablof{U02}     & 0.0188 & 0.9987 & 5.3809 & 0.82 \\
\ablof{U03}     & 0.0202 & 0.9984 & 5.3496 & 0.63 \\
\ablof{U04}     & 0.0297 & 0.9976 & 5.4055 & 0.30 \\
\ablof{U05}     & 0.0168 & 0.9987 & 5.3222 & 0.36 \\
\ablof{U06}     & 0.0208 & 0.9985 & 5.3406 & 0.62 \\
\ablof{U07}     & 0.0173 & 0.9986 & 5.3162 & 0.40 \\
\ablof{U08}     & 0.0076 & 0.9994 & 5.2800 & 0.55 \\
\ablof{U10}     & 0.0202 & 0.9984 & 5.3494 & 0.62 \\
\ablof{U11}     & 0.0202 & 0.9984 & 5.3494 & 0.63 \\
\bottomrule
\end{tabular}\end{table}

\subsection{Median Performance --- Group Ablations}
\label{sec:results_group}

\begin{table}[H]
\centering
\caption{Median performance --- group ablations.}
\label{tab:res_group}
\renewcommand{\arraystretch}{1.25}
\setlength{\tabcolsep}{6pt}
\small
\begin{tabular}{l l
    S[table-format=1.4]
    S[table-format=1.4]
    S[table-format=1.4]
    S[table-format=2.2]}
\toprule
\textbf{Variant} & \textbf{Units removed} &
  {\textbf{IGD}$^+$ $\downarrow$} &
  {\textbf{nHV} $\uparrow$} &
  {\textbf{SP} $\downarrow$} &
  {\textbf{Time (s)} $\downarrow$} \\
\midrule
G-SEARCH  & U01 + U02         & 0.0222 & 0.9980 & 5.3734 & 0.60 \\
G-ADAPT   & U05 + U06 + U07   & 0.0193 & 0.9985 & 5.3165 & 0.35 \\
G-ARCHIVE & U08 + U10         & 0.0076 & 0.9994 & 5.2800 & 0.54 \\
G-NOADV   & U08 + U10 + U11   & 0.0076 & 0.9994 & 5.2800 & 0.54 \\
G-NOBASE  & U01--U07          & 0.0514 & 0.9940 & 5.4332 & 0.28 \\
G-CORE    & U02,U03,U08,U10,U11 & 0.0075 & 0.9995 & 5.2847 & 0.55 \\
\bottomrule
\end{tabular}
\end{table}

\subsection{Robustness: Inter-Quartile Range}
\label{sec:results_iqr}

Table~\ref{tab:iqr} reports the inter-quartile range (IQR) of IGD$^+$
and nHV to characterise result robustness across the 30 independent runs.
Narrower IQR indicates more consistent behaviour.

\begin{table}[H]
\centering
\caption{Robustness --- IQR of IGD$^+$ and nHV over 30 runs.}
\label{tab:iqr}
\renewcommand{\arraystretch}{1.2}
\setlength{\tabcolsep}{8pt}
\small
\begin{tabular}{l
    S[table-format=1.4] S[table-format=1.4]
    S[table-format=1.4] S[table-format=1.4]}
\toprule
 & \multicolumn{2}{c}{\textbf{IGD}$^+$}
 & \multicolumn{2}{c}{\textbf{nHV}} \\
\cmidrule(lr){2-3}\cmidrule(lr){4-5}
\textbf{Variant} & {Median} & {IQR} & {Median} & {IQR} \\
\midrule
\mosho{} (full) & 0.0202 & 0.0033 & 0.9984 & 0.0004 \\
MOSHO-Base      & 0.0249 & 0.0031 & 0.9983 & 0.0005 \\
\midrule
\ablof{U01}     & 0.0206 & 0.0039 & 0.9981 & 0.0006 \\
\ablof{U02}     & 0.0188 & 0.0034 & 0.9987 & 0.0004 \\
\ablof{U03}     & 0.0202 & 0.0029 & 0.9984 & 0.0004 \\
\ablof{U04}     & 0.0297 & 0.0046 & 0.9976 & 0.0007 \\
\ablof{U05}     & 0.0168 & 0.0032 & 0.9987 & 0.0005 \\
\ablof{U06}     & 0.0208 & 0.0030 & 0.9985 & 0.0002 \\
\ablof{U07}     & 0.0173 & 0.0049 & 0.9986 & 0.0006 \\
\ablof{U08}     & 0.0076 & 0.0020 & 0.9994 & 0.0002 \\
\ablof{U10}     & 0.0202 & 0.0033 & 0.9984 & 0.0004 \\
\ablof{U11}     & 0.0202 & 0.0033 & 0.9984 & 0.0004 \\
\midrule
G-SEARCH        & 0.0222 & 0.0049 & 0.9980 & 0.0009 \\
G-ADAPT         & 0.0193 & 0.0027 & 0.9985 & 0.0004 \\
G-ARCHIVE       & 0.0076 & 0.0020 & 0.9994 & 0.0002 \\
G-NOADV         & 0.0076 & 0.0020 & 0.9994 & 0.0002 \\
G-NOBASE        & 0.0514 & 0.0062 & 0.9940 & 0.0010 \\
G-CORE          & 0.0075 & 0.0020 & 0.9995 & 0.0002 \\
\bottomrule
\end{tabular}\end{table}

% ============================================================
\section{Performance Impact Analysis}
\label{sec:impact}
% ============================================================

\subsection{Normalised Impact Scores}
\label{sec:impact_norm}

The \emph{normalised impact score} for unit $U_i$ on metric $m$ is:

\begin{equation}
  \mathrm{Impact}_{\mathrm{nHV}}(U_i) =
  \frac{\mathrm{nHV}(\text{\mosho{}}) - \mathrm{nHV}(\text{\ablof{$U_i$}})}
       {\mathrm{nHV}(\text{\mosho{}}) - \mathrm{nHV}(\text{MOSHO-Base})}
  \times 100\%
  \label{eq:impact}
\end{equation}

\noindent For IGD$^+$, the numerator and denominator are reversed because
lower is better:
$(\mathrm{IGD}^+(\text{\ablof{$U_i$}})-\mathrm{IGD}^+(\text{\mosho{}})) /
(\mathrm{IGD}^+(\text{MOSHO-Base})-\mathrm{IGD}^+(\text{\mosho{}}))$.
A score of 100\% means the unit explains the full gap between full
\mosho{} and MOSHO-Base.  Negative scores mean that removing the unit
improved the corresponding median metric in this execution.

\begin{table}[H]
\centering
\caption{Normalised impact scores (\%) on nHV and IGD$^+$.
  Colour bands:
  \colorbox{critred!25}{$>$60\%~critical},
  \colorbox{imptorange!25}{40--60\%~important},
  \colorbox{modblue!20}{20--40\%~moderate},
  \colorbox{marggreen!20}{$<$20\%~marginal}.}
\label{tab:impact}
\renewcommand{\arraystretch}{1.25}
\setlength{\tabcolsep}{10pt}
\begin{tabular}{l c c c}
\toprule
\textbf{Unit} & \textbf{Name} &
  \textbf{Impact on nHV (\%)} &
  \textbf{Impact on IGD}$^+$ \textbf{(\%)} \\
\midrule
U04 & Explicit acceptance  & \cellcolor{critred!25}515.1 & \cellcolor{critred!25}202.8 \\
U01 & Discrete kernels     & \cellcolor{critred!25}199.0 & \cellcolor{marggreen!20}~9.3 \\
U06 & Dual adaptive signal & \cellcolor{marggreen!20}-32.0 & \cellcolor{marggreen!20}14.1 \\
U03 & Deferred truncation  & \cellcolor{marggreen!20}10.5 & \cellcolor{marggreen!20}~0.0 \\
U10 & Extreme solutions    & \cellcolor{marggreen!20}~0.0 & \cellcolor{marggreen!20}~0.0 \\
U11 & Mono-obj.\ fallback  & \cellcolor{marggreen!20}~0.0 & \cellcolor{marggreen!20}~0.0 \\
U07 & Active archive       & \cellcolor{marggreen!20}-89.6 & \cellcolor{marggreen!20}-60.8 \\
U02 & Uniform crossover    & \cellcolor{marggreen!20}-138.6 & \cellcolor{marggreen!20}-29.6 \\
U05 & Energy + restart     & \cellcolor{marggreen!20}-150.2 & \cellcolor{marggreen!20}-71.8 \\
U08 & K-means archive      & \cellcolor{marggreen!20}-604.5 & \cellcolor{marggreen!20}-266.6 \\
\bottomrule
\end{tabular}
\end{table}

\subsection{Absolute Performance Deltas}
\label{sec:impact_delta}

Table~\ref{tab:delta} reports the absolute performance delta against
full \mosho{}.  For lower-is-better metrics (IGD$^+$, SP, time),
$\Delta = \text{variant} - \text{\mosho{}}$; for nHV,
$\Delta = \text{\mosho{}} - \text{variant}$.  Positive quality deltas
therefore indicate that \mosho{} is better; negative time deltas indicate
that the variant is faster.

\begin{table}[H]
\centering
\caption{Absolute performance deltas.
  Positive $\Delta$IGD$^+$ and $\Delta$SP indicate \mosho{} is better;
  positive $\Delta$nHV indicates \mosho{} achieves higher hypervolume;
  negative $\Delta$Time indicates the variant is faster.}
\label{tab:delta}
\renewcommand{\arraystretch}{1.2}
\setlength{\tabcolsep}{7pt}
\small
\begin{tabular}{l
    S[table-format=+1.4]
    S[table-format=+1.4]
    S[table-format=+1.4]
    S[table-format=+2.2]}
\toprule
\textbf{Variant} &
  {$\Delta$\,IGD$^+$} &
  {$\Delta$\,nHV} &
  {$\Delta$\,SP} &
  {$\Delta$\,Time (s)} \\
\midrule
MOSHO-Base      & +0.0047 & +0.0002 & +0.0570 & -0.03 \\
\mowso{}        & \multicolumn{4}{c}{not run in \texttt{results\_ablation\_study}} \\
\midrule
\ablof{U01}     & +0.0004 & +0.0003 & +0.0142 & -0.04 \\
\ablof{U02}     & -0.0014 & -0.0002 & +0.0314 & +0.17 \\
\ablof{U03}     & +0.0000 & +0.0000 & +0.0001 & -0.02 \\
\ablof{U04}     & +0.0095 & +0.0008 & +0.0560 & -0.35 \\
\ablof{U05}     & -0.0034 & -0.0002 & -0.0272 & -0.30 \\
\ablof{U06}     & +0.0007 & -0.0001 & -0.0088 & -0.03 \\
\ablof{U07}     & -0.0029 & -0.0001 & -0.0332 & -0.25 \\
\ablof{U08}     & -0.0125 & -0.0010 & -0.0695 & -0.10 \\
\ablof{U10}     & +0.0000 & +0.0000 & +0.0000 & -0.03 \\
\ablof{U11}     & +0.0000 & +0.0000 & +0.0000 & -0.02 \\
\bottomrule
\end{tabular}\end{table}

\noindent\textbf{Key observation on execution time.}
The fastest OAT variant is \ablof{U04} (0.30~s), reflecting the cost
saved by unconditional acceptance.  \ablof{U02} is the only OAT variant
slower than full \mosho{} in median runtime (+0.17~s), while U10 and U11
are numerically identical in quality to full \mosho{} and slightly faster.

% ============================================================
\section{Statistical Analysis}
\label{sec:stat}
% ============================================================

\subsection{Statistical Pipeline}
\label{sec:stat_pipeline}

The statistical framework adopted here is a \textbf{three-step mandatory
suite}, complemented by an optional fourth step for added robustness.
Each step is motivated by the design and the conventions of the
evolutionary computation literature.

\begin{table}[H]
\centering
\caption{Statistical analysis pipeline.}
\label{tab:pipeline}
\renewcommand{\arraystretch}{1.35}
\setlength{\tabcolsep}{6pt}
\small
\begin{tabular}{c l p{6.5cm} l}
\toprule
\textbf{Step} & \textbf{Method} & \textbf{Purpose} & \textbf{Status} \\
\midrule
1 & Wilcoxon Signed-Rank Test (paired)
  & Test whether removing each unit $U_i$ significantly degrades
    performance.  Directly matches the paired-seed design and is
    distribution-free.
  & \textbf{\textcolor{pipelinered}{Mandatory}} \\
2 & Holm--Bonferroni Correction
  & Control the Family-Wise Error Rate across $m=10$ simultaneous
    comparisons. 
  & \textbf{\textcolor{pipelineorange}{Mandatory}} \\
3 & Vargha--Delaney $A_{12}$ Effect Size
  & Quantify the \emph{magnitude} of each unit's contribution beyond
    binary significance; required for tier classification.
  & \textbf{\textcolor{pipelineyellow}{Mandatory}} \\
4 & Bootstrap 95\% Confidence Intervals
  & Provide interval estimates for median differences and $A_{12}$;
    strengthens robustness claims.
  & \textcolor{pipelinegreen}{\textbf{Optional}} \\
\bottomrule
\end{tabular}
\end{table}


\subsection{Step 1 --- Wilcoxon Signed-Rank Test}
\label{sec:stat_wilcoxon}

When per-unit ablation runs are available, the \textbf{Wilcoxon
signed-rank test} compares \mosho{} and \ablof{$U_i$} over the $R=30$
paired runs:

\[
d_r = \mathrm{IGD}^+(\ablof{$U_i$},\,r)
      - \mathrm{IGD}^+(\mosho{},\,r),
\quad r = 1,\dots,30.
\]

\noindent\textbf{Hypotheses:}
\begin{itemize}[leftmargin=2em, itemsep=1pt]
  \item $H_0$: the median of $d_r$ is zero (no difference).
  \item $H_1$: the median of $d_r$ is greater than zero
        (removal increases IGD$^+$, i.e.\ degrades performance).
\end{itemize}

Significance level: $\alpha = 0.05$.  A significant positive shift
confirms a beneficial contribution of the removed unit.

\subsection{Step 2 --- Multiple Comparison Control: Holm--Bonferroni}
\label{sec:stat_holm}

With $m = 10$ simultaneous tests, unadjusted $p$-values inflate the
probability of false positives.  The \textbf{Holm--Bonferroni procedure}
controls the Family-Wise Error Rate (FWER) while retaining greater
statistical power than classical Bonferroni.

Let $p_1 \leq p_2 \leq \dots \leq p_m$ be the sorted $p$-values.
Each is compared against the adjusted threshold:

\[
\alpha_i = \frac{0.05}{m - i + 1}.
\]

The procedure rejects $H_0$ for all $p_j$ with $j \leq i^*$, where
$i^*$ is the last index satisfying $p_{i^*} \leq \alpha_{i^*}$.
All subsequent hypotheses are retained.

\subsection{Step 3 --- Effect Size: Vargha--Delaney $A_{12}$}
\label{sec:stat_a12}

Statistical significance alone does not convey practical importance.
The \textbf{Vargha--Delaney $A_{12}$} effect size is therefore computed
for every comparison:

\[
A_{12}(\mosho{},\,\ablof{$U_i$})
=
P\!\bigl(\mosho{} < \ablof{$U_i$}\bigr)
+ \tfrac{1}{2}\,
P\!\bigl(\mosho{} = \ablof{$U_i$}\bigr).
\]

This is the probability that \mosho{} achieves a lower (better) IGD$^+$
than the ablated variant.  Interpretation thresholds follow
Table~\ref{tab:a12_thresholds}.

\begin{table}[H]
\centering
\caption{Vargha--Delaney $A_{12}$ interpretation thresholds.}
\label{tab:a12_thresholds}
\renewcommand{\arraystretch}{1.3}
\small
\begin{tabular}{cll}
\toprule
$A_{12}$ & \textbf{Magnitude} & \textbf{Interpretation} \\
\midrule
0.50           & None   & No observable difference \\
0.51 -- 0.56   & Small  & Weak advantage for \mosho{} \\
0.57 -- 0.71   & Medium & Moderate advantage for \mosho{} \\
$> 0.71$       & Large  & Strong advantage for \mosho{} \\
\bottomrule
\end{tabular}
\end{table}

\subsection{Decision Rule}
\label{sec:stat_decision}

Each unit is classified by combining adjusted $p$-value and $A_{12}$:

\begin{itemize}[leftmargin=2em, itemsep=4pt]
  \item Adjusted $p < 0.05$ \textbf{and} $A_{12} > 0.71$:
        \quad\critical{Critical contribution.}
  \item Adjusted $p < 0.05$ \textbf{and} $0.57 \leq A_{12} \leq 0.71$:
        \quad\important{Important contribution.}
  \item Adjusted $p < 0.05$ \textbf{and} $0.51 \leq A_{12} < 0.57$:
        \quad\secondary{Moderate contribution.}
  \item Adjusted $p \geq 0.05$:
        \quad\neutral{No statistically significant quality contribution.}
\end{itemize}

\noindent This combined criterion ensures conclusions are both
statistically valid and practically meaningful.

\subsection{Results}
\label{sec:stat_results}

\subsubsection{Pairwise Test Results (Wilcoxon + $A_{12}$)}
\label{sec:stat_pairwise}

Table~\ref{tab:wilcoxon} reports the Wilcoxon signed-rank test results
for the OAT ablations available in the supplied result root.  The same
paired-run aggregation is used as in the descriptive tables.

\begin{table}[H]
\centering
\caption{Wilcoxon signed-rank results on IGD$^+$ and nHV
  (30 paired runs, Holm--Bonferroni adjusted).
  $H_0$: no difference vs.\ full \mosho{}.
  Significance: $^{*}p<0.05$;\ $^{**}p<0.01$;
  $^{***}p<0.001$;\ \textsf{ns}~not significant.}
\label{tab:wilcoxon}
\renewcommand{\arraystretch}{1.2}
\setlength{\tabcolsep}{5pt}
\small
\begin{tabular}{l
    S[table-format=1.4] c c
    S[table-format=1.4] c c}
\toprule
 & \multicolumn{3}{c}{\textbf{IGD}$^+$}
 & \multicolumn{3}{c}{\textbf{nHV}} \\
\cmidrule(lr){2-4}\cmidrule(lr){5-7}
\textbf{Variant} &
  {$p$-value} & {Sig.} & {$\hat{A}_{12}$} &
  {$p$-value} & {Sig.} & {$\hat{A}_{12}$} \\
\midrule
\ablof{U01} & 0.0000 & $^{***}$ & 0.62 & 0.0000 & $^{***}$ & 0.73 \\
\ablof{U02} & 1.0000 & \textsf{ns} & 0.37 & 1.0000 & \textsf{ns} & 0.40 \\
\ablof{U03} & 0.0000 & $^{***}$ & 0.50 & 0.0000 & $^{***}$ & 0.50 \\
\ablof{U04} & 0.0000 & $^{***}$ & 1.00 & 0.0000 & $^{***}$ & 0.98 \\
\ablof{U05} & 1.0000 & \textsf{ns} & 0.16 & 1.0000 & \textsf{ns} & 0.31 \\
\ablof{U06} & 0.0000 & $^{***}$ & 0.61 & 1.0000 & \textsf{ns} & 0.43 \\
\ablof{U07} & 1.0000 & \textsf{ns} & 0.18 & 1.0000 & \textsf{ns} & 0.45 \\
\ablof{U08} & 1.0000 & \textsf{ns} & 0.00 & 1.0000 & \textsf{ns} & 0.00 \\
\ablof{U10} & 1.0000 & \textsf{ns} & 0.50 & 1.0000 & \textsf{ns} & 0.50 \\
\ablof{U11} & 1.0000 & \textsf{ns} & 0.50 & 1.0000 & \textsf{ns} & 0.50 \\
\bottomrule
\end{tabular}
\end{table}

\subsubsection{Unit Ranking by Median IGD$^+$}
\label{sec:stat_ranking}

Table~\ref{tab:ranking} ranks the reference and OAT ablation variants by
median IGD$^+$ and cross-references the $A_{12}$ effect size against full
\mosho{} where applicable.

\begin{table}[H]
\centering
\caption{Variant ranking by median IGD$^+$ (ascending = better),
  cross-referenced with $\hat{A}_{12}$ against full \mosho{}.}
\label{tab:ranking}
\renewcommand{\arraystretch}{1.2}
\setlength{\tabcolsep}{9pt}
\small
\begin{tabular}{l c c c c}
\toprule
\textbf{Variant} &
  \textbf{Median IGD}$^+$ &
  \textbf{Rank} &
  $\hat{A}_{12}$ &
  \textbf{Interpretation} \\
\midrule
\ablof{U08}     & 0.0076 & 1  & 0.00 & Ablation improves median IGD$^+$ \\
\ablof{U05}     & 0.0168 & 2  & 0.16 & Ablation improves median IGD$^+$ \\
\ablof{U07}     & 0.0173 & 3  & 0.18 & Ablation improves median IGD$^+$ \\
\ablof{U02}     & 0.0188 & 4  & 0.37 & Ablation improves median IGD$^+$ \\
\ablof{U10}     & 0.0202 & 5  & 0.50 & No observable median change \\
\ablof{U11}     & 0.0202 & 6  & 0.50 & No observable median change \\
\mosho{} (full) & 0.0202 & 7  & ---  & Reference \\
\ablof{U03}     & 0.0202 & 8  & 0.50 & No observable median change \\
\ablof{U01}     & 0.0206 & 9  & 0.62 & Full \mosho{} improves quality \\
\ablof{U06}     & 0.0208 & 10 & 0.61 & Full \mosho{} improves quality \\
MOSHO-Base      & 0.0249 & 11 & ---  & Reference baseline \\
\ablof{U04}     & 0.0297 & 12 & 1.00 & Full \mosho{} improves quality \\
\bottomrule
\end{tabular}
\end{table}

\subsection{Step 4 --- Bootstrap Confidence Intervals (Optional)}
\label{sec:stat_bootstrap}

Bootstrap 95\% confidence intervals may be computed for:
\begin{itemize}[leftmargin=2em, itemsep=2pt]
  \item the median difference in IGD$^+$
        ($\Delta_{\mathrm{med}} = \mathrm{med}(d_r)$), and
  \item the Vargha--Delaney $A_{12}$ effect size.
\end{itemize}

Resampling is performed with replacement over the $R=30$ paired runs
($10^4$ resamples).  A confidence interval for $\Delta_{\mathrm{med}}$
that excludes zero, or for $A_{12}$ that excludes 0.5, reinforces the
Wilcoxon conclusions and provides the most rigorous evidence available
for the contribution of each unit.  This step is recommended for
camera-ready or journal submissions.

% ============================================================
\section{Interpretation and Discussion}
\label{sec:interp}
% ============================================================

\subsection{Unit Classification Summary}
\label{sec:interp_class}

Table~\ref{tab:classification} consolidates the current OAT evidence.
Here, positive impact means that removing the unit degraded the median
relative to full \mosho{}; negative impact means the ablation improved
the median in this execution.

\begin{table}[H]
\centering
\caption{Final unit classification across all evidence sources.}
\label{tab:classification}
\renewcommand{\arraystretch}{1.3}
\setlength{\tabcolsep}{5pt}
\small
\begin{tabular}{c l c c c l}
\toprule
\textbf{Unit} & \textbf{Name} &
  \textbf{Impact (\%)} & $\hat{A}_{12}$ &
  \textbf{Wilcoxon} & \textbf{Tier} \\
\midrule
U04 & Explicit acceptance  & 515.1 & 1.00 & $^{***}$ & \textbf{Critical} \\
U01 & Discrete kernels     & 199.0 & 0.62 & $^{***}$ & \textbf{Important} \\
U06 & Dual adaptive signal & -32.0 & 0.61 & $^{***}$ & \textbf{Important on IGD$^+$ only} \\
U03 & Deferred truncation  & 10.5 & 0.50 & $^{***}$ & \textbf{Statistical, negligible effect} \\
U10 & Extreme solutions    & 0.0 & 0.50 & \textsf{ns} & \textbf{No observed effect} \\
U11 & Mono-obj.\ fallback  & 0.0 & 0.50 & \textsf{ns} & \textbf{No observed effect} \\
U02 & Uniform crossover    & -138.6 & 0.37 & \textsf{ns} & \textbf{Ablation improves median} \\
U05 & Energy + restart     & -150.2 & 0.16 & \textsf{ns} & \textbf{Ablation improves median} \\
U07 & Active archive       & -89.6 & 0.18 & \textsf{ns} & \textbf{Ablation improves median} \\
U08 & K-means archive      & -604.5 & 0.00 & \textsf{ns} & \textbf{Ablation improves median} \\
\bottomrule
\multicolumn{6}{l}{Impact column is nHV impact normalised by the MOSHO vs MOSHO-Base gap.} \\
\end{tabular}
\end{table}

\subsection{Key Findings}
\label{sec:interp_findings}

\begin{enumerate}[leftmargin=*, nosep, itemsep=6pt]

  \item \textbf{Full \mosho{} improves over MOSHO-Base, but the gain is
  small.}
  The median IGD$^+$ improves from 0.0249 to 0.0202 and nHV improves
  from 0.9983 to 0.9984.  MOSHO-Base is slightly faster in median
  runtime (0.62~s versus 0.65~s for full \mosho{}).

  \item \textbf{Explicit acceptance (U04) is the clearest positive unit.}
  Removing U04 gives the worst OAT median IGD$^+$ (0.0297) and nHV
  (0.9976), with large $A_{12}$ values and Holm-adjusted significance.

  \item \textbf{Several ablations outperform full \mosho{} on median
  quality.}
  Removing U08, U05, U07, and U02 improves median IGD$^+$ relative to
  full \mosho{}.  U08 is especially notable: \ablof{U08}, G-ARCHIVE, and
  G-NOADV all obtain IGD$^+\approx0.0076$ and nHV$\approx0.9994$.

  \item \textbf{U10 and U11 have no measurable effect in this execution.}
  \ablof{U10} and \ablof{U11} are numerically identical to full \mosho{}
  on the reported quality medians, which is consistent with the code path:
  U08 cluster sampling is inactive for the $5^6$ search space, and U11
  only activates when the latency-degeneracy condition is met.

\end{enumerate}


\section{Conclusion}
\end{document}
